% ============================================================================
% Cụm tuyên bố cuối bài: Lời cảm ơn + Đạo đức + Mã nguồn — đều ngắn gọn (10pt).
% ============================================================================

\section*{Lời cảm ơn}
\label{sec:ack}

{\fontsize{10}{12}\selectfont
Tác giả cảm ơn PGS.\,TS.\,Phạm Thọ Hoàn (người hướng dẫn chính) và Công ty Cổ phần Khoa học và Công nghệ Giáo dục OLM đã hỗ trợ nghiên cứu này.
\par}

\section*{Đạo đức và dữ liệu}

{\fontsize{10}{12}\selectfont
\noindent Dữ liệu do OLM cung cấp theo điều khoản dịch vụ, phục vụ mục đích nghiên cứu khoa học giáo dục phi thương mại. Định danh học sinh đã được hash hoá trước khi phân tích, đảm bảo không thể nhận diện ngược cá nhân. Bài báo chỉ công bố các kết quả tổng hợp; dữ liệu chi tiết của từng học sinh được tác giả chia sẻ riêng cho mục đích tái lập nghiên cứu khi có yêu cầu chính đáng từ các nhà nghiên cứu khác.
\par}

\section*{Khai báo sử dụng trí tuệ nhân tạo}

{\fontsize{10}{12}\selectfont
\noindent Trong quá trình thực hiện nghiên cứu, tác giả sử dụng mô hình ngôn ngữ lớn Claude (Anthropic), thông qua công cụ Claude Code, vào hai việc: (i) hỗ trợ viết và gỡ lỗi mã nguồn phân tích bằng R và Python của quy trình B0--B7; (ii) hỗ trợ tìm kiếm và tóm tắt tài liệu liên quan. Toàn bộ mã nguồn do AI hỗ trợ đã được tác giả rà soát và chạy lại trên dữ liệu; mọi tài liệu trích dẫn đã được tác giả đối chiếu với bản gốc. Tác giả chịu hoàn toàn trách nhiệm về nội dung, số liệu và kết luận của bài báo.
\par}

\section*{Mã nguồn}

{\fontsize{10}{12}\selectfont
\noindent Mã nguồn quy trình phân tích và các kết quả tổng hợp có tại \url{https://github.com/dtnntd/irt-lsem}.
\par}
